O pipeline experimental permitiu avaliar filtragem FIR, descritores espectrais
e respostas Gabor sob a restrição de treino apenas com batimentos normais. A
execução inter-paciente na MIT-BIH mostrou que o acréscimo de processamento e
features não produz ganho monotônico. One-Class SVM favoreceu a representação
bruta, Isolation Forest respondeu melhor aos descritores espectrais, e LOF teve
seu melhor resultado com features Gabor.

Esses resultados preservam a hipótese como questão experimental, não como
conclusão fechada. A próxima etapa deve concentrar-se em análise de limiar por
FPR alvo, seleção ou compressão de features Gabor, comparação com baselines
supervisionados apenas como referência externa, e inspeção de exemplos de erro.
O enquadramento deve permanecer como benchmark de DSP e detecção one-class, sem
promessa de diagnóstico clínico.
